%
% teil3.tex -- Beispiel-File für Teil 3
%
% (c) 2020 Prof Dr Andreas Müller, Hochschule Rapperswil
%
% !TEX root = ../../buch.tex
% !TEX encoding = UTF-8
%
\section{Analytische Fortsetzung
\label{Analytische Fortsetzung}}

Die analytische Fortsetzung ist ein Konzept aus der komplexen Analysis.
Im Grunde genommen kann durch sie eine Funktion erweitert werden
ausserhalb ihres Definitionsbereiches.

Voraussetzung dafür ist lediglich, dass die Funktion holomorph ist.

Grundsätzlich wird in der Analytischen Fortsetzung einfach eine neue Funktion konstruiert, die den selben Definitionsbereich wie die ursprüngliche Funktion hat. 
Jedoch ist diese neue Funktion auch ausserhalb dieses Bereichs definiert.

In der Anwendung auf die Riemannsche Zeta-Funktion kann eine solche analytische Fortsetzung fast auf die gesamte komplexe Ebene angewendet werden. 
Die einzige Ausnahme ist der Punkt bei s = 1.

%
%

\subsection{Anwendung auf die Zeta-Funktion\label{Anwendung auf die Zeta-Funktion}}

Beispielsweise wenn man die Zeta-Funktion anschaut mit

\[\zeta(s) = 1 + \frac{1}{2^{s}} + \frac{1}{3^{s}} + \frac{1}{4^{s}} + \ldots\]

, so ist diese nur korrekt für alle Zahlen \textgreater{} 1.

Wenn man jedoch dann die Analytische Fortsetzung verwendet und die Funktion anders darstellt durch eine Erweiterung des Terms, bekommt man

\[\zeta(s) = \frac{1}{{1 - 2}^{1 - s}}(1 + \frac{1}{2^{s}} + \frac{1}{3^{s}} + \frac{1}{4^{s}} + \ldots),\]

was nun schon korrekt ist für alle Zahlen \textgreater{} 0.

Die Formel wird nun erweitert, bis man

\[\zeta(s) = 2^s \pi^{s-1} \sin\left(\frac{\pi s}{2}\right) \Gamma(1 - s)\,\zeta(1 - s)\].

hat. 
So ist sie definiert für auch fast alle negativen Werte für s. 
Wenn man für s jetzt -1 einsetzt, so kommt man mit den Umformungen von

\[\zeta(-1) = 2^{-1} \pi^{-2} \sin\left(-\frac{\pi}{2}\right)\Gamma(2)\zeta(2)\].

auf

\[\sin\left(-\frac{\pi}{2}\right) = -1, \quad \Gamma(2) = 1! = 1, \quad \zeta(2) = \frac{\pi^2}{6}\].

und schlussendlich auf

\[
\frac{1}{2} \cdot \frac{1}{\pi^2} \cdot (-1) \cdot 1 \cdot \frac{\pi^2}{6} = -\frac{1}{12},
\]
was
\[
-\frac{1}{12}
\] ergibt. 

%Berechnung und analytische Fortsetzung noch besser beschreiben, auch Zusammenhang mit komplexen Zahlen näher erklären.


